\documentclass[11pt,a4paper]{article}
\usepackage[T1]{fontenc}\usepackage[utf8]{inputenc}\usepackage{lmodern}\usepackage{microtype}
\usepackage[a4paper,margin=2.25cm]{geometry}\usepackage{longtable,booktabs,array,calc}\usepackage{fancyvrb}
\usepackage{xcolor}\usepackage{hyperref}\usepackage{xurl}\usepackage{fancyhdr}\usepackage{enumitem}\usepackage{parskip}\usepackage{titlesec}
\definecolor{ddtadark}{RGB}{30,38,48}\definecolor{ddtablue}{RGB}{38,82,126}
\hypersetup{colorlinks=true,linkcolor=ddtablue,urlcolor=ddtablue,citecolor=ddtablue}
\pagestyle{fancy}\fancyhf{}\lhead{Documentation-Driven Threat Analysis}\rhead{BA3 work plan R14}\cfoot{\thepage}\setlength{\headheight}{14pt}
\setlist[itemize]{leftmargin=1.5em,itemsep=0.2em,topsep=0.25em}\setlist[enumerate]{leftmargin=1.7em,itemsep=0.2em,topsep=0.25em}
\titleformat{\section}{\Large\bfseries\color{ddtadark}}{}{0em}{}\titleformat{\subsection}{\large\bfseries\color{ddtadark}}{}{0em}{}\titleformat{\subsubsection}{\normalsize\bfseries\color{ddtadark}}{}{0em}{}
\providecommand{\tightlist}{\setlength{\itemsep}{0pt}\setlength{\parskip}{0pt}}\setlength{\emergencystretch}{4em}\hbadness=10000\hfuzz=20pt\sloppy
\begin{document}
\section{DDTA research work plan after documentation-layer
closure}\label{ddta-research-work-plan-after-documentation-layer-closure}

\textbf{WORK PLAN - REVISION 14}

\textbf{Status:} active BA3 execution plan after BA3-T2 cross-baseline
identity/lifecycle pressure testing.

\textbf{Supersedes:} Revision 13 only for forward execution state;
R1-R13 remain historical research records.

\subsection{1. Fixed prior state}\label{fixed-prior-state}

\begin{itemize}
\tightlist
\item
  Chapters 2-4: CLOSED / FINAL for current thesis scope.
\item
  Documentation layer: CLOSED.
\item
  BA0-R systems-modeling prior-art research: CLOSED.
\item
  BA0 responsibility/non-goals: CLOSED.
\item
  BA1 minimal BAE identity ontology: CLOSED.
\item
  BA2 relation/action vocabulary: CLOSED BY BA2-T4.
\item
  BA3-T1 provenance/origin lower bound: COMPLETED / PROVISIONAL PASS.
\item
  \texttt{BAReferent} and \texttt{BAProposition}: ACCEPTED first-class
  semantic identity families.
\item
  ThreatForge is a case study/tooling instrument, not DDTA semantic
  authority.
\end{itemize}

\subsection{2. BA3 objective}\label{ba3-objective}

BA3 defines the smallest derivation/provenance/authority mechanics that
connect governed documentation to Base Analysis while preserving project
authority in governed documentation and making change/revalidation
reproducible.

\subsection{3. BA3-T1 - provenance/origin lower
bound}\label{ba3-t1---provenanceorigin-lower-bound}

\textbf{Status: COMPLETED / PROVISIONAL PASS WITH LOWER-BOUND
CANDIDATE.}

T1 requires independent provenance on both BA1 identity families,
many-to-many source lineage, immutable baseline context, exact source
locator, the three origin states and explicit derivation basis/rationale
for derived meaning.

\subsection{4. BA3-T2 - cross-baseline identity, staleness and
lifecycle}\label{ba3-t2---cross-baseline-identity-staleness-and-lifecycle}

\textbf{Status: COMPLETED / PROVISIONAL PASS WITH IDENTITY-LIFECYCLE
REFINEMENT.}

T2 replayed the facial-access M1/M2/M3 mutations and the
order-fulfillment internal/external inventory-authority counterfactual.

\subsubsection{Main dispositions}\label{main-dispositions}

\begin{verbatim}
BAReferent continuity          retain when reusable referent meaning survives
BAProposition continuity       retain only under normalized assertion equivalence
source revision == BA revise  REJECTED
BA continuity dispositions     RETAIN | REPLACE | RETIRE
BA-level REVISE                NOT REQUIRED
review state                   PENDING_REVIEW | ACCEPTED | REJECTED
freshness                      CURRENT | STALE
accepted REPLACE               prior element SUPERSEDED
accepted RETIRE                prior element RETIRED
resolved diagnostic            RETIRE or REPLACE, never mutate origin to GROUNDED
BA1 / BA2 reopen               NOT TRIGGERED
\end{verbatim}

\subsubsection{Mutation summary}\label{mutation-summary}

\begin{itemize}
\tightlist
\item
  M1 Ethernet -\textgreater{} Wi-Fi retains abstract connectivity but
  replaces the realization proposition.
\item
  M2 external -\textgreater{} project-owned transport replaces
  responsibility polarity, may retire external consumption and can
  introduce a new transport branch without changing the consumer
  transfer/security meaning.
\item
  M3 remote -\textgreater{} local recognition retires the remote
  transfer and its old local constraints rather than dragging them onto
  the new behavior.
\item
  external/internal WMS pressure produces the same
  responsibility-boundary pattern in a distinct domain.
\end{itemize}

\subsection{5. Why BA3 remains open}\label{why-ba3-remains-open}

T2 still does not close:

\begin{itemize}
\tightlist
\item
  exact derivation-rule/rationale registry and reproducibility;
\item
  effective governed context/dependency semantics sufficient for
  principled staleness propagation;
\item
  source-to-analysis and analysis-to-source feedback lineage after
  accepted correction;
\item
  full BA3 closure review.
\end{itemize}

These must be tested before BA4 begins.

\subsection{6. BA3-T3 - derivation-rule reproducibility and
change-impact lineage pressure
test}\label{ba3-t3---derivation-rule-reproducibility-and-change-impact-lineage-pressure-test}

\textbf{NEXT / NOT EXECUTED BY THIS PLAN REVISION.}

BA3-T3 must attack whether the T1/T2 contracts are sufficient for
reproducible derivation and localized change impact without inventing a
general dependency graph.

It must pressure-test at least:

\begin{enumerate}
\def\labelenumi{\arabic{enumi}.}
\tightlist
\item
  the facial M4 raw-capture -\textgreater{} opaque-reference mutation
  where sibling governed context changes FR-3.4 meaning;
\item
  provider-specific raw state -\textgreater{} governed result
  normalization in order/payment/WMS cases;
\item
  the minimum identity and content required of
  \texttt{derivationRuleRef} or an equivalent inspectable rationale;
\item
  whether derived elements can be rebuilt/reviewed deterministically
  enough from \texttt{derivationBasis\ +\ rule/rationale};
\item
  how direct source changes, derivation-basis changes and participant
  replacement mark only potentially affected BA elements stale;
\item
  whether explicit effective-context/dependency links are required
  beyond source parentage;
\item
  how an analysis-produced corrective candidate points back to the
  governed source area without becoming source authority;
\item
  how an accepted governed correction reconnects to the next BA baseline
  without turning analysis output into project truth;
\item
  whether any result forces BA1/BA2 or the BA3-T1/T2 lower bound to
  reopen.
\end{enumerate}

\subsubsection{BA3-T3 guardrails}\label{ba3-t3-guardrails}

Do not:

\begin{itemize}
\tightlist
\item
  design BA4 projections;
\item
  define formal STRIDE/STRIDE-AI overlay schemas;
\item
  define AnalysisRecord/Common Finding;
\item
  design ThreatForge implementation classes/database tables;
\item
  turn \texttt{derivationRuleRef} into an unrestricted programming
  language;
\item
  make arbitrary legacy-document extraction a requirement.
\end{itemize}

\subsubsection{BA3-T3 exit condition}\label{ba3-t3-exit-condition}

Produce a falsifiable derivation/change-impact candidate and identify
whether a distinct BA3 closure review is still required. Do not close
BA3 merely because one implementation strategy can encode the examples.

\subsection{7. BA4 - projections}\label{ba4---projections}

Only after BA3 closes, define derived human and method projections from
the same accepted BA semantics. No view becomes project authority.

\subsection{8. BA5 - lexical vocabulary and optional
assistance}\label{ba5---lexical-vocabulary-and-optional-assistance}

Only after BA3/BA4 boundaries are stable, define display labels,
authoring synonyms and optional assistance. NLP/LLM proposals remain
candidate-producing support.

\subsection{9. BA6 - complete Base Analysis regression and
closure}\label{ba6---complete-base-analysis-regression-and-closure}

Regress the complete BA0-BA5 design against the closed corpora and at
least one structurally different holdout. BA6 may reopen only the
smallest earlier responsibility on a material counterexample.

\subsection{10. Analysis envelope remains
downstream}\label{analysis-envelope-remains-downstream}

Only after Base Analysis closure:

\begin{itemize}
\tightlist
\item
  A1 - AnalysisRecord / execution envelope;
\item
  A2 - common reviewed Finding boundary;
\item
  A3 - change/provenance integration;
\item
  formal methodology overlays and STRIDE/STRIDE-AI evaluations.
\end{itemize}

\subsection{11. Next authorized
microstep}\label{next-authorized-microstep}

Only after the BA3-T2 package is reviewed, committed, pushed and
remotely verified, execute:

\begin{quote}
\textbf{BA3-T3 - derivation-rule reproducibility and change-impact
lineage pressure test.}
\end{quote}

Do not start BA4, formal threat-method overlays, Common Finding schema
or implementation work.

\end{document}
